% ==========================================
% TASK 3
% ==========================================
\section{Task 3: SPP with Prefetch-Aware Dynamic I-RRIP}\label{sec:task3}

Task~3 runs the final Task~1 prefetcher (SPP, \texttt{prefetcher/spp.l1d\_pref}, unchanged) at the L1D
together with the dynamic I-RRIP policy of Task~2B at the L2, and makes that policy aware of the
lines SPP brings in. All prefetch handling lives in \texttt{replacement/dyn\_irrip.l2c\_repl} behind one
switch, \texttt{PF\_AWARE}. With no prefetcher none of the new code runs, so the Task~2B results are
unchanged.

\subsection{What the L2 sees from SPP}\label{sec:task3-l2view}

Every SPP prefetch is issued at the L1D and, when it misses there, travels to the L2 as a
\texttt{PREFETCH} request. SPP chooses its fill level from the L1D MSHR occupancy (Task~1):
\begin{itemize}[nosep]
  \item \textbf{FILL\_L1} (L1D MSHR less than half full): the line is filled into the L2 \emph{and} the L1D.
        The demand that follows hits in the L1D, so the L2 copy is a duplicate that the L2 will
        rarely be asked for again.
  \item \textbf{FILL\_L2} (L1D MSHR at least half full): the line is filled into the L2 \emph{only}.
        It is not in the L1D, so its demand access has to come to the L2.
\end{itemize}
The two kinds of prefetch therefore have very different value \emph{to the L2}, even though both
arrive with the same packet type. We measured this directly (\cref{tab:task3-firstuse},
\cref{fig:t3-firstuse}): only 0.15--0.72\% of FILL\_L1 prefetch fills ever get a demand hit in the L2,
while 13--19\% of FILL\_L2 prefetch fills do. FILL\_L2 prefetches are only 1.5--4\% of the prefetch fills,
yet they account for about half or more of all L2 demand hits on prefetched lines.

\begin{table}[htbp]
\centering
\caption{Prefetch fills into the L2 and their first demand use, split by fill level
(SPP + prefetch-aware Dyn.\ I-RRIP, whole run including warm-up).}\label{tab:task3-firstuse}
\tablesetup
\begin{tabular}{lrrrrrr}
\toprule
 & \multicolumn{3}{c}{\colhead{FILL\_L1 (L1D + L2)}} & \multicolumn{3}{c}{\colhead{FILL\_L2 (L2 only)}} \\
\cmidrule(lr){2-4}\cmidrule(lr){5-7}
\colhead{Trace} & \colhead{Fills} & \colhead{First use} & \colhead{Used (\%)} & \colhead{Fills} & \colhead{First use} & \colhead{Used (\%)} \\
\midrule
Trace 1 & 143403 & 1028 & 0.72 & 5455 & 1062 & 19.47 \\
Trace 2 &  94376 &  232 & 0.25 & 1421 &  230 & 16.19 \\
Trace 3 & 106602 &  158 & 0.15 & 3220 &  423 & 13.14 \\
\bottomrule
\end{tabular}
\end{table}

\resultfig[0.72\linewidth]{task3_pf_first_use}{Share of prefetched L2 lines that receive a demand
hit in the L2, by fill level. L2-only prefetches are 25--90$\times$ more likely to be used in the L2.}{fig:t3-firstuse}

\subsection{Design: fill-level-aware prefetch handling}\label{sec:task3-design}

The L2 fill path in \texttt{src/cache.cc} stores the fill level of the MSHR entry in a new
\texttt{CACHE} member, \texttt{repl\_fill\_level}, just before it calls
\texttt{update\_replacement\_state()} (one line in \texttt{cache.cc}, one in \texttt{inc/cache.h}).
The replacement interface is otherwise unchanged. The policy then applies these rules in both
follower modes (I-RRIP and RRIP) and in the leader sets:

\begin{enumerate}[nosep]
  \item \textbf{FILL\_L1 prefetch fill} $\rightarrow$ RRPV 7 (\texttt{PF\_L1\_INSERT}), the most distant
        position. The line is a duplicate of an L1D copy and is the first candidate for eviction, so it
        cannot push out instruction lines or demand data.
  \item \textbf{FILL\_L2 prefetch fill} $\rightarrow$ inserted like a demand data miss of the current mode:
        RRPV 4 in I-RRIP mode, RRPV 6 in RRIP mode. The line has to survive in the L2 until its demand arrives.
  \item \textbf{Prefetch hit in the L2} $\rightarrow$ RRPV unchanged. The request only copies the line up
        to the L1D, which says nothing about future L2 reuse (the Task~2B policy decremented the RRPV here).
  \item \textbf{First demand hit on a prefetched line} $\rightarrow$ promoted by the normal LOAD-hit rule of
        the mode (to RRPV $\le 2$ for data in I-RRIP mode, to RRPV 4 in RRIP mode). The prefetch has
        proved useful, so the line is treated like any reused line.
  \item \textbf{Set dueling sees demand traffic only.} \texttt{PREFETCH} requests are not counted in the
        leader-set access/miss counters or in the epoch length. Prefetch traffic depends on the
        prefetcher, not on how well the policy keeps instruction or data lines, so counting it would
        dilute the instruction ratio that the duel uses to choose the mode.
\end{enumerate}
No per-line state is added for the policy itself: the RRPV array and the existing
\texttt{BLOCK::\allowbreak prefetch} bit are enough (a 2-bit per-line tag is kept only for the statistics of
\cref{tab:task3-firstuse}). \Cref{tab:task3-policy} summarises how each class of line is treated.

\begin{table}[htbp]
\centering
\caption{Treatment of each line class (RRPV range 0--7) in the prefetch-aware policy.}\label{tab:task3-policy}
\tablesetup
\begin{tabularx}{\linewidth}{lLLL}
\toprule
\colhead{Line class} & \colhead{Insertion (I-RRIP / RRIP mode)} & \colhead{Promotion on hit} & \colhead{Eviction priority} \\
\midrule
Instruction (demand)     & 1 / 6 & to 0 / to 4 & last: among RRPV-7 lines, data is evicted first in I-RRIP mode \\
Data (demand)            & 4 / 6 & to $\le 2$ / to 4 & normal \\
Prefetch, FILL\_L1       & 7 / 7 & prefetch hit: none; first demand hit: as data & first (inserted at the eviction position) \\
Prefetch, FILL\_L2       & 4 / 6 & prefetch hit: none; first demand hit: as data & as demand data \\
\bottomrule
\end{tabularx}
\end{table}

\subsection{Rationale for Handling Prefetched Lines}\label{sec:task3-rationale}

\begin{itemize}
  \item \textbf{Most prefetched L2 lines are dead on arrival.} SPP is an L1D prefetcher: a useful prefetch
        is consumed in the L1D, and the L2 copy is not asked for again unless the L1D evicts it first.
        Under 1\% of FILL\_L1 prefetch fills are ever used in the L2 (\cref{tab:task3-firstuse}).
        Inserting them at the eviction position means that inaccurate or duplicate prefetches never
        displace the instruction lines that I-RRIP is protecting.
  \item \textbf{The useful L2 prefetches must not be thrown away.} A FILL\_L2 prefetch is the only copy
        above the LLC. Inserting it at the distant position would evict it before its demand arrives and
        waste the DRAM access. Giving it the same position as demand data keeps it as long as any
        other data line. The 25--90$\times$ difference in L2 use between the two kinds of prefetch is the
        measured reason for treating them differently instead of treating ``prefetch'' as one class.
  \item \textbf{A prefetch hit is not reuse.} A prefetch request that hits in the L2 is about to copy
        the line into the L1D; promoting the L2 copy would keep a line that the L1D now serves.
        Only a demand hit is evidence of L2 reuse.
  \item \textbf{The duel must judge the policy, not the prefetcher.} Prefetch misses are data misses
        that both leader groups would suffer whatever the mode. Keeping them out of the counters makes
        the I-RRIP/RRIP decision depend only on demand behaviour, the same signal as in Task~2B.
\end{itemize}

\subsection{Results}\label{sec:task3-results}

We compare four configurations, all built from the same code base and run with 25M warm-up and 25M
simulated instructions: the baseline (no prefetcher, LRU at the L2), SPP with LRU (Task~1 final),
dynamic I-RRIP without a prefetcher (Task~2B), and SPP with the prefetch-aware dynamic I-RRIP
(Task~3, written SPP + Dyn.\ I-RRIP-PF in the tables). Speedup is over the baseline. Accuracy and coverage are defined as in Task~1
(coverage = baseline L1D load misses removed). Means are geometric for IPC and speedup and
arithmetic otherwise.

\begin{table}[htbp]
\centering
\caption{Task 3: IPC, speedup and L1D prefetch metrics.}\label{tab:task3-perf}
\tablesetup
\begin{tabular}{llrrrr}
\toprule
\colhead{Trace} & \colhead{Configuration} & \colhead{IPC} & \colhead{Speedup} & \colhead{Accuracy (\%)} & \colhead{Coverage (\%)} \\
\midrule
\multirow{4}{*}{Trace 1} & Baseline (no pf, LRU)    & 0.2320 & 1.0000 & --   & --    \\
                         & SPP + LRU                & 0.2339 & 1.0080 & 47.6 & 10.87 \\
                         & Dyn.\ I-RRIP (no pf)     & 0.2320 & 1.0000 & --   & --    \\
                         & SPP + Dyn.\ I-RRIP-PF & 0.2338 & 1.0077 & 47.4 & 10.90 \\
\midrule
\multirow{4}{*}{Trace 2} & Baseline (no pf, LRU)    & 0.2506 & 1.0000 & --   & --    \\
                         & SPP + LRU                & 0.2541 & 1.0141 & 51.6 & 4.64  \\
                         & Dyn.\ I-RRIP (no pf)     & 0.2547 & 1.0165 & --   & --    \\
                         & SPP + Dyn.\ I-RRIP-PF & \textbf{0.2586} & \textbf{1.0322} & 52.0 & 4.61 \\
\midrule
\multirow{4}{*}{Trace 3} & Baseline (no pf, LRU)    & 0.3819 & 1.0000 & --   & --    \\
                         & SPP + LRU                & 0.3882 & 1.0165 & 33.4 & 15.28 \\
                         & Dyn.\ I-RRIP (no pf)     & 0.3817 & 0.9996 & --   & --    \\
                         & SPP + Dyn.\ I-RRIP-PF & 0.3880 & 1.0161 & 33.4 & 15.37 \\
\midrule
\multirow{4}{*}{Mean}    & Baseline (no pf, LRU)    & 0.2811 & 1.0000 & --   & --    \\
                         & SPP + LRU                & 0.2847 & 1.0129 & 44.2 & 10.26 \\
                         & Dyn.\ I-RRIP (no pf)     & 0.2826 & 1.0053 & --   & --    \\
                         & SPP + Dyn.\ I-RRIP-PF & \textbf{0.2863} & \textbf{1.0186} & 44.3 & 10.29 \\
\bottomrule
\end{tabular}
\end{table}

\begin{table}[htbp]
\centering
\caption{Task 3: load MPKI of the L1D, the L2C, and the L2 split into instruction and data loads.}\label{tab:task3-mpki}
\tablesetup
\begin{tabular}{llrrrr}
\toprule
\colhead{Trace} & \colhead{Configuration} & \colhead{L1D load} & \colhead{L2C load} & \colhead{L2 instr.} & \colhead{L2 data} \\
\midrule
\multirow{4}{*}{Trace 1} & Baseline (no pf, LRU)    & 10.68 & 17.72 & 10.37 & 7.35 \\
                         & SPP + LRU                &  9.52 & 16.75 & 10.41 & 6.35 \\
                         & Dyn.\ I-RRIP (no pf)     & 10.68 & 17.63 &  9.97 & 7.67 \\
                         & SPP + Dyn.\ I-RRIP-PF &  9.51 & 16.52 &  9.94 & 6.59 \\
\midrule
\multirow{4}{*}{Trace 2} & Baseline (no pf, LRU)    & 13.54 & 28.19 & 19.58 & 8.61 \\
                         & SPP + LRU                & 12.91 & 27.70 & 19.66 & 8.04 \\
                         & Dyn.\ I-RRIP (no pf)     & 13.54 & 24.16 & 15.10 & 9.05 \\
                         & SPP + Dyn.\ I-RRIP-PF & 12.91 & 23.40 & 14.95 & 8.44 \\
\midrule
\multirow{4}{*}{Trace 3} & Baseline (no pf, LRU)    &  3.40 &  3.96 &  1.13 & 2.84 \\
                         & SPP + LRU                &  2.88 &  3.34 &  1.16 & 2.18 \\
                         & Dyn.\ I-RRIP (no pf)     &  3.40 &  3.98 &  1.10 & 2.87 \\
                         & SPP + Dyn.\ I-RRIP-PF &  2.88 &  3.32 &  1.11 & 2.20 \\
\midrule
\multirow{4}{*}{Mean}    & Baseline (no pf, LRU)    &  9.21 & 16.63 & 10.36 & 6.26 \\
                         & SPP + LRU                &  8.44 & 15.93 & 10.41 & 5.52 \\
                         & Dyn.\ I-RRIP (no pf)     &  9.21 & 15.26 &  8.72 & 6.53 \\
                         & SPP + Dyn.\ I-RRIP-PF &  \textbf{8.44} & \textbf{14.41} &  \textbf{8.67} & 5.75 \\
\bottomrule
\end{tabular}
\end{table}

\resultfig{task3_ipc}{Task 3: IPC of each configuration.}{fig:t3-ipc}
\resultfig{task3_speedup}{Task 3: speedup over the baseline (no prefetcher, LRU). Bar labels are the speedup.}{fig:t3-speedup}
\resultfig[0.95\linewidth]{task3_mpki}{Task 3: load MPKI of the L1D, the L2C, L2 instruction and L2 data loads.}{fig:t3-mpki}

\begin{figure}[htbp]
\centering
\IfFileExists{figures/task3_accuracy.png}{\includegraphics[width=0.49\linewidth]{figures/task3_accuracy.png}}{}%
\hfill
\IfFileExists{figures/task3_coverage.png}{\includegraphics[width=0.49\linewidth]{figures/task3_coverage.png}}{}
\caption{Task 3: L1D prefetch accuracy (left) and coverage (right) of SPP with LRU and with the
prefetch-aware Dyn.\ I-RRIP at the L2.}\label{fig:t3-acccov}
\end{figure}

\subsection{Discussion}\label{sec:task3-discussion}

\paragraph{Outcome.} SPP with the prefetch-aware dynamic I-RRIP reaches a geometric-mean speedup of
\textbf{1.0186} (+1.86\%), against 1.0129 for SPP alone and 1.0053 for dynamic I-RRIP alone, so it
outperforms both on average. It beats dynamic I-RRIP on every trace.

\paragraph{Why the gains add up.} The two mechanisms remove different misses. SPP removes L1D data
misses (mean L1D load MPKI 9.21 $\rightarrow$ 8.44, coverage about 10\%), which an L2 policy cannot touch.
Dynamic I-RRIP removes L2 instruction misses (mean 10.36 $\rightarrow$ 8.67), which SPP cannot touch.
On Trace~2, the most instruction-bound trace, SPP gives +1.41\% and dynamic I-RRIP +1.65\%; together
they give +3.22\%, slightly more than the product of the two (1.0308). The L2 instruction MPKI drops from
19.66 to 14.95 while the L1D load MPKI stays at the SPP level (12.91).

\paragraph{SPP is not disturbed.} Accuracy and coverage are the same with LRU and with the new policy
(mean accuracy 44.2\% vs 44.3\%, coverage 10.26\% vs 10.29\%). The L2 policy changes how long the L2 keeps
lines, not what SPP predicts. Because FILL\_L1 prefetches are inserted at the eviction position, the
extra prefetch traffic does not cost I-RRIP its instruction-miss savings: the L2 instruction MPKI with
SPP (8.67) is as low as without a prefetcher (8.72).

\paragraph{Trace 1 and Trace 3.} On these traces dynamic I-RRIP alone gains nothing (1.0000 and 0.9996):
Trace~3 has almost no L2 instruction misses (1.13 MPKI), and on Trace~1 the instruction misses saved
(10.37 $\rightarrow$ 9.97) are paid back in data misses (7.35 $\rightarrow$ 7.67). With no instruction-miss
saving to add, the combination can only match SPP, and it does so to within 0.04\%
(1.0077 vs 1.0080 and 1.0161 vs 1.0165, a difference of about 30k cycles in 107M on Trace~1).
The set-dueling threshold of Task~2B still selects I-RRIP for most epochs on these traces
(97\% on Trace~1, 75\% on Trace~3); a stricter mode decision is the natural next step for these
workloads.

\paragraph{Cost of distant insertion.} Inserting FILL\_L1 prefetches at RRPV 7 lowers the L2 hit rate of
\emph{prefetch requests}: when SPP re-prefetches a line that the L1D has already evicted, the L2 copy has
often gone too (Trace~1: prefetch-request hit rate 27\% with LRU vs 18\%). These re-prefetches go to
the LLC instead. The effect on IPC is small, since the demand hit rate of the L2 rises at the same
time (Trace~1 L2 load MPKI 16.75 $\rightarrow$ 16.52).
